\documentclass{article}
\usepackage{../../style/header}

\begin{document}

\title{Algebraic Topology}
\author{Lectured by Steven Sivek \\
Scribed by Yu Coughlin}
\date{Spring 2026}

\maketitle

\tableofcontents

\newpage

\section{Covering spaces}

\begin{definition}
    A space $X$ is \textbf{locally connected} if for $x\in X$ and $x\in V\subseteq X$ open, there exists an open $x\in U\subset V$ which is connected, we can do the same for path-connected.
\end{definition}

For this entire section, all spaces are assumed to be path-connected and locally path-connected.

\begin{definition}
    A continuous map $p:\tilde{X}\rightarrow X$ is \textbf{covering map} if for all $x\in X$, there exists an open neighborhood $x\in U\subset X$ such that \[
    p^{-1}(U) = \bigsqcup_{\alpha}V_\alpha
    \]is a nonemtpy disjoint union of open $V_\alpha\subset \tilde{X}$, with $p\vert_{V_\alpha}:V_\alpha\rightarrow U$ a homeomorphism for all $\alpha$.
\end{definition}

\begin{lemma}
    Covering maps are surjective.
\end{lemma}

\subsection{Lifting}

\begin{lemma}
    If $p:\tilde{X}\rightarrow X$ is a covering space, with $\tilde{f}_0,\tilde{f}_1:Y\rightarrow \tilde{X}$ two lifts of $f:Y\rightarrow X$, then the set \[
    S = \{y\in Y \mid \tilde{f}_0(y) = \tilde{f}_1(y)
    \] is both open and closed in $Y$.
    \begin{proof}
        For all $y\in S$, there exists an open neighbourhood $f(y)\in U \subset X$ with $p^{-1}(U) = \bigsqcup_\alpha V_\alpha$, then $\tilde{f}_0(y) = \tilde{f}_1(y) \in V_\beta$, and thus on $N:=\tilde{f}_0^{-1}(V_\beta)\cap\tilde{f}_1^{-1}(V_\beta)$ we have \[
        p\vert_{V_\beta}\circ \tilde{f}_0\vert_N = f\vert_N = p\vert_{V_\beta}\circ \tilde{f}_1\vert_N
        \] but $p\vert_{V_\beta}$ is a homeomorphism, so $\tilde{f}_0\vert_N = \tilde{f}_1\vert_N$, with $y\in N\subset S$ open, thus $S$ is open.

        For all $y\in \overline{S}\setminus S$, there exists an open neighbourhood $f(y)\in U\subset X$ with $p^{-1}(U) = \bigsqcup_\alpha V_\alpha$, and thus $\tilde{f}_0(y)\in V_\beta$ and $\tilde{f}_1(y) \in V_\gamma$ for $\beta\neq\gamma$, so the open neighbourhood $\tilde{f}_0^{-1}(V_\beta)\cap\tilde{f}_1^{-1}(V_\gamma)$ of $y$ intersects $S$, but then $V_\beta\cap V_\gamma\neq \emptyset$, a contradiction.
    \end{proof}
\end{lemma}

\begin{theorem}[Homotopy lifting theorem]
    For $p:\tilde{X}\rightarrow X$ a covering map, and $Y$ a locally connected space, for all homotopies $H:Y\times I \rightarrow X$ with a lift $\tilde{h}_0:Y\times\{0\}\rightarrow\tilde{X}$ of $h_0(y) = H(Y,0)$, there there exists a unique homotopy \[
    \begin{tikzcd}
        Y\times\{0\} \arrow[r, "h_0"] \arrow[d, hook] & \tilde{X} \arrow[d, "p"] \\
        Y\times I \arrow[r, "H"] \arrow[ur, "\exists ! \tilde{H}"] & X
    \end{tikzcd}
    \]\begin{proof}
        For each $y\in Y$, consider the map $\gamma_y:I\rightarrow X$ by $\gamma_y(t) = H(y,t)$, if we have two lifts $\tilde{H},\tilde{K}$ of $H$, then they agree at $(y,0)$, and thus by the previous lemma, as $I$ is connected, they agree on all of $I$, hence $\tilde{H}=\tilde{K}$.

        Now, fix $y_0\in Y$, for all $t\in I$, there exists open $y_0\in V_t\subset Y$ and $t \in W_t\subset I$, such that $H(V_t\times \overline{W}_t)$ lies in a single evenly covered $U_\alpha$, as $I$ is compact, we can cover this with finitely many $W_t$ and let $V$ be the intersection of the corresponding $V_t$. As $Y$ is locally connected, we may shrink $V$ to be connected and consider the $H(V\times[t_i,t_{i+1}])$ each contained in some $U_{\alpha_i}$ for $i = 0,1,\ldots,n-1$; now, starting with $\tilde{h}_0$, we can extend each $\tilde{H}(V\times[0,t_i])$ to $\tilde{H}(V\times [0,t_{i+1}])$ by taking the unique lift of the connected $V\times [0,t_{i+1}]$ lying interely in the same $V_{\beta_i}$ over $U_{\alpha_i}$ as $V\times \{t_i\}$.

        We now have, for all $y\in Y$, a lift $\tilde{H}_y:V_y\times I\rightarrow \tilde{X}$, but by the previous lemma, these all agree on all their intersections ,so we can glue these all together to $\tilde{H}:Y\times I \rightarrow \tilde{X}$.
    \end{proof}
\end{theorem}

\begin{corollary}[Path lifting theorem]
    For $p:\tilde{X}\rightarrow X$ a covering map, and $\gamma:I\rightarrow X$ a path, then for each $\tilde{x}\in p^{-1}(\gamma(0))$, there exists a unique path $\tilde{\gamma}:I\rightarrow \tilde{X}$ lifting $\gamma$ to $\tilde{X}$ and starting at $\tilde{x}$.
    \begin{proof}
        This is the homotopy lifting theorem with $Y = *$.
    \end{proof}
\end{corollary}

\subsection{Covers and the fundamental group}

\begin{definition}
    For $(X,x)$ a pointed topological space, let $\mathbf{Cov}(X,x)$ be the full subcategory of $\mathbf{Top}_*/(X,x)$ of covering maps. We get the functor $\pi_1:\mathbf{Cov}(X,x)\rightarrow\mathbf{Grp}$.
\end{definition}

\begin{proposition}
    If $p:(Y,y)\rightarrow (X,x)$ is a covering map, then $p_*:\pi_1(Y,y)\rightarrow \pi_1(X,x)$ is injective.\begin{proof}
        If $\gamma:I\rightarrow Y$ has $p_*([\gamma]) = [e_{x}]$, in $X$, then the homotopy $p\circ\gamma\simeq e_x$ has a unique lift to $\gamma\simeq e_y$.
    \end{proof}
\end{proposition}

So can associate to each cover $p:Y\rightarrow X$ a subgroup $p_*\pi_1(Y,y)\leq \pi_1(X,x)$.

\begin{theorem}[Lifting criterion]
    If $p:(Y,y)\rightarrow (X,x)$ is a covering map and $f:(Z,z)\rightarrow (X,x)$ continuous with $Z$ path-connected and locally path-connected, then we can lift $f$ iff $f_*\pi_1(Z,z)\subseteq p_*\pi_1(Y,y)$.
    \begin{proof}
        The ($\Rightarrow$) direction is obvious, so we suppose $f_*\pi_1(Z,z)\subseteq p_*\pi_1(Y,y)$, and define the lift $\tilde{f}$. For all $w\in Z$, there exists a path $\gamma:I\rightarrow Z$ from $z$ to $w$, we can lift $f\circ \gamma$ to some unique $g$ with $g(0) = y$, now take $\tilde{f}(w):= g(1)$.

        First, to show $\tilde{f}$ is well-defined, suppose we have another path $\gamma':I\rightarrow Z$ from $z$ to $w$, then consider $\phi = f\circ(\gamma * \overline{\gamma}')\in f_*\pi_1(Z,z)\subseteq p_*\pi_1(Y,y)$, so we can find a loop $\psi:I\rightarrow Y$ based at $y$ with $p_*[\psi] = [\phi]$, so there is a homotopy $H$ from $p\circ\psi$ to $\phi$ in $X$, which uniquely lifts to some $\tilde{H}$ from $\psi$ to some lift $\tilde{\phi}$ of $\phi$, which must be a loop based at $y$. But, we can form another lift of $\phi$ by following $g$, and then $g'$, a lift of $\overline{f\circ\gamma'}$ starting at $g(1)$, by uniquness, this lift must be a loop at $y$, so the unique lift of $f\circ\gamma'$ must start at $y$ and end of $g(1)$, so $\tilde{f}(w)$ doesn't depend on our choice of path.

        Now, to show $\tilde{f}$ is continuous, let $U\subset X$ be an open, path-connected, evenly covered neighbourhood of $f(w)$, with $V_\alpha$ the unique sheet containing $\tilde{f}(w)$, as the $U,V$ form a basis of $X,Y$ respectively, it suffices to show $\tilde{f}^{-1}(V)$ contains an open neighbourhood of $w$. Let $W\subset f^{-1}(U)\subset Z$ be a path-connected open neighbourhood of $w$, then for all $w'\in W$, there exists a path $\eta:I\rightarrow W$ from $w$ to $w'$, its image $f\circ\eta$ lies entirely within $U$, where $p$ is a homeomorphism from $V$, so we have the lift $p^{-1}\circ f \circ \eta$ whose starting point is $\tilde{f}(w)$, by the uniqueness of lifts, the endpoint of $p^{-1}\circ f \circ \eta \in V$ is $\tilde{f}(w')$, therefore $W\subset \tilde{f}^{-1}(V)$, hence $\tilde{f}$ is continuous.
    \end{proof}
\end{theorem}

\begin{corollary}
    If $p:(Y,y)\rightarrow (X,x)$ and  $q:(Z,z)\rightarrow (X,x)$ are two non-isomorphic coverings of $(X,x)$, then $p_*\pi_1(Y,y)\neq q_*\pi_1(Z,z)$.\begin{proof}
        Suppose $p_*\pi_1(Y,y)=q_*\pi_1(Z,z)$, then the lifting criterion gives $f:(Y,y)\rightarrow (Z,z)$ and $g:(Z,z)\rightarrow (Y,y)$, but now $f\circ g$ and $\id_Y$ are both lifts of $p$ along $p$, and so by uniqueness are equal.
    \end{proof}
\end{corollary}

\begin{proposition}
    If in the statement of the lifting criterion, $f=q$ is a covering map, then the lift $r$ is a covering map.
    \begin{proof}
        Take $v\in Y$, then $p(v)\in X$ has evenly open neighbourhoods $U_y$ and $U_z$ by $p$ and $q$ respectively, take $U\subset U_y\cap U_z$ connected, and have its sheets be $V_\alpha$ and $W_\beta$ in $p$ and $q$ respectively, let $V_0$ be the sheet containing $v$. For all $r(w)\in V_0$, with $w\in Z$, we have $q(w)=pr(w)\in U$, so $w\in W_\beta$ for some $\beta$, as $W_\beta$ is a connected, and $r(W_\beta)$ is a subset of $p^{-1}(U)$, we must have $r(W_\beta)\subset V_0$, with $r\vert_{W_\beta} = (p\vert_{V_0})^{-1}\circ q\vert_{W_\beta}$ a homeomorphism, so $r^{-1}(V_0)$ is a disjoint union of some of the $W_\beta$, with each $r\vert_{W_\beta}$ a homeomorphism onto $V_0$, hence $r$ is a covering map.
    \end{proof}
\end{proposition}

\subsection{Deck transformations}

\begin{theorem}
    If $p:Y\rightarrow X$ is a covering map, then for each $x\in X$, there is a right action $p^{-1}(x)\curvearrowleft \pi_1(X,x)$. For all $y\in p^{-1}(x)$ and $\gamma:I\rightarrow X$, let $\tilde{\gamma}$ be the unique lift of $\gamma$ to $Y$ with starting point $y$, then take $y\cdot [\gamma] = \tilde{\gamma}(1)\in p^{-1}(x)$.
    \begin{proof}
        To show this is well-defined, suppose $\gamma,\gamma'$ are loops at $x$ in $X$, with $H$ a homotopy from $\gamma$ to $\gamma'$, then there is a unique lift \[
        \begin{tikzcd}
            I\times \{0\} \arrow[r, "\tilde{\gamma}"] \arrow[d, hook] & Y \arrow[d, "p"] \\
            I\times I \arrow[ur, dashed, "\tilde{H}"] \arrow[r, "H"] & X
        \end{tikzcd}
        \] which must satisfy $\tilde{H}(0,t) = \tilde{\gamma}(0) = y$ and $\tilde{H}(1,t) = \tilde{\gamma}(1)$ for all $t$, which implies $\tilde{\gamma}'(1)=\tilde{\gamma}(1)$, so $y\cdot[\gamma]=y\cdot[\gamma']$.

        To check this is a group homomorphism, we have to show $(y\cdot [\gamma])\cdot[\gamma'] = y\cdot([\gamma]\cdot[\gamma'])$, let $\tilde{\gamma}''$ be the unique lift of $\gamma*\gamma'$ starting at $y$, then the RHS is $\tilde{\gamma}''(1)$. But the first half of $\tilde{\gamma}''$ must be the unique lift of $\gamma$ starting at $y$, and the second half the unique lift of $\gamma'$ starting and $\tilde{\gamma}(1)$, so $\tilde{\gamma}''(1)=\tilde{\gamma}'(1)$.
    \end{proof}
\end{theorem}

As $Y$ is path-connected, this action is clearly transitive.

\begin{proposition}
    For $p:Y\rightarrow X$ a covering with $x\in X$, the stabiliser of $y\in p^{-1}(x)$, is $p_*\pi(Y,y)$.
    \begin{proof}
        We can lift any loop $\gamma$ in the stabiliser to a loop $\tilde{\gamma}$ in $Y$, then $\gamma=p_*([\tilde{\gamma}])$.
    \end{proof}
\end{proposition}

By orbit-stabiliser, we define the \textbf{degree} $\deg(p) :=\abs{p^{-1}(x)} = [\pi_1(X,x):p_*\pi_1(Y,y)]$.

\begin{proposition}
    For $p:Y\rightarrow X$ a covering with $x\in X$ and $y\in p^{-1}(x)$, then for all $[\gamma]\in\pi_1(X,x)$, we have $p_*\pi_1(Y,y\cdot[\gamma]) = [\gamma]^{-1}p_*\pi_1(Y,y)[\gamma]$.
    \begin{proof}
        Let $\tilde{\gamma}$ be the lift of $\gamma$ from $y$ to $z$ in $Y$, this induces an isomorphism $\phi_{\tilde{\gamma}}:\pi_1(Y,y)\rightarrow\pi_1(Y,z)$ by $[\psi]\mapsto[\overline{\tilde{\gamma}}*\psi*\tilde{\gamma}]$, and so $p_*\phi_{\tilde{\gamma}}([\psi]) = [(p\circ\overline{\tilde{\gamma}})*(p\circ\psi)*(p\circ\tilde{\gamma})]= [\overline{\gamma}*(p\circ \psi) * \gamma] = [\gamma]^{-1}p_*([\psi])[\gamma]$.
    \end{proof}
\end{proposition}

\begin{definition}
    We call an isomorphism of $p:(Y,y)\rightarrow (X,x)$ in $\mathbf{Cov}(X,x)$ a \textbf{Deck transformation} and write $\Deck(p)$ or $\Deck(Y)$ when clear from context. If $\Deck(Y)$ acts transitively on $p^{-1}(x)$, then we call $p$ \textbf{normal}.
\end{definition}

\begin{proposition}
    For $p:Y\rightarrow X$ a covering, with $x\in X$, and $y,z\in p^{-1}(x)$, then there exists $f\in\Deck(Y)$ with $f(y)=z$ iff $p_*\pi_1(Y,y)=p_*\pi_1(Y,z)$.
    \begin{proposition}
        Such a deck transformation is exactly a lift of $p$ against $p$, so this is just the lifting criterion.
    \end{proposition}
\end{proposition}

\begin{theorem}
    A covering $p:(Y,y)\rightarrow (X,x)$ is regular iff $p_*\pi_1(Y,y)\unlhd \pi_1(X,x)$.\begin{proof}
        This is combining all the previous props.
    \end{proof}
\end{theorem}

\begin{theorem}
    For $p:(Y,y)\rightarrow (X,x)$ a covering, $\Deck(Y)\cong N_{\pi_1(X,x)}(p_*\pi_1(X,x))/p_*\pi_1(X,x)$\begin{proof}
        Write $G = \pi_1(X,x)$ and $H=p_*\pi_1(Y,y)$, and define a map $\phi:N_G(H)\rightarrow \Deck(Y)$ by for all $[\gamma]\in N_G(H)$, as $p_*\pi_1(Y,y\cdot[\gamma]) = [\gamma]^{-1}p_*\pi_1(Y,y)[\gamma] = p_*\pi_1(Y,y)$ which gives a funique deck transformation $f=\phi([\gamma])$ such that $f(y)=y\cdot[\gamma]$.

        To show $\phi$ is a homomorphism, take $[\gamma],[\gamma']\in N_G(H)$ with $\phi([\gamma])=f$ and $\phi([\gamma'])=g$, then we can lift $\gamma'$ to a path $\alpha$ from $y\cdot [\gamma]$ to $(y\cdot [\gamma])\cdot [\gamma']$, then $f^{-1}\circ \alpha$ is the unique lift of $\gamma'$ starting at $y$, so $f^{-1}(\alpha(1)) = g(y)$, thus $\phi[\gamma]\circ\phi[\gamma'](y) = f(g(y)) = \alpha(1)= \phi[\gamma *\gamma'](y)$.

        For surjectivity, $f\in \Deck(Y)$ sending $y$ to $z\in p^{-1}(x)$, then the trasitive $\pi_1(X,x)$ action on $p^{-1}(x)$ gives $\gamma$ with $z=y\cdot[\gamma]$, so $p_*\pi_1(X,y)=p_*\pi_1(X,y\cdot [\gamma])$, hence $\gamma$ is in the normaliser.

        $[\gamma]\in N_G(H)$ has $\phi[\gamma]=\id_Y$ iff $y\cdot[\gamma]=y$ iff $[\gamma]\in P_*\pi_1(Y,y)=H$.
    \end{proof}
\end{theorem}

\subsection{Universal covers}

The lifting criterion for covering spaces says is $(\tilde{X},\tilde{x})$ is simply connected, then it is initial in $\mathbf{Cov}(X,x)$, the previous theorem thus says $\pi_1(X,x)\cong \Deck(\tilde{X})$, we want to know when these exists.

\begin{definition}
    A space $X$ is \textbf{semi-locally simply connected} if all $x\in X$ has an open neighbourhood $U\subset X$ such that the induced map from the inclusion $i_*\pi_1(U,x)\rightarrow \pi_1(X,x)$ is trivial.
\end{definition}

This is weaker than locally simply connected as we allow loops in our neighbourhood to be trivial only when considered in all of $\pi_1(X,x)$. Consider the cone on the Hawaiian earing for something semi-locally simply connected but no locally simply-connected

\begin{theorem}
    If $X$ is path-connected, locally path-connected, and semi-locally simply connected, then $X$ admites a simply connected cover $p:(\tilde{X},\tilde{x}_0)\rightarrow (X,x_0)$.\begin{proof}
        First consider the set $\tilde{X} = \{\gamma:I\rightarrow X \mid \gamma(0) = x\}$ module homotopy rel the endpoints, and let $\tilde{x}=[e_x]$ and $p:\tilde{X}\rightarrow X$ send $[\gamma]$ to $\gamma(1)$, wich is clearly well-defined.

        To define a topology on this set, first consider the set $\cU$ of path-connected open sets $U\subset X$ where $i_*:\pi_1(U,z)\rightarrow \pi_1(X,z)$ is trivial for some $z\in U$, as $X$ is locally path-connected and semi-locally simply connected, this is a basis for the topology on $X$. Now define a basis on $\tilde{X}$, for $U\in \cU$ and $\gamma:I\rightarrow X$ from $x_0$ to some $x\in U$, consider $V_{U,[\gamma]} = \{[\gamma * \eta] \mid \eta :I\rightarrow U, \ \eta(0)=\gamma(1)\}$. We can immediately check: if $[\gamma] = [\gamma']$ then $V_{U,[\gamma]} = V_{U,[\gamma']}$, if $[\gamma']\in V_{U,[\gamma]}$ then $V_{U,[\gamma]}=V_{U,[\gamma']}$, and if $\gamma(1)\in W \in \cU$, then $V_{W,[\gamma]}\subset V_{U,[\gamma]}$, so let these form a basis for the topology on $\tilde{X}$.
    \end{proof}
    TODO: finish the proof
\end{theorem}

\subsection{Fundamental theorem}

\begin{theorem}
    For $(X,x)$ path-connected, locally path-connected, and semi-locally simply connected, there is a bijection between isomorphism class of covers of $(X,x)$ and subgroup of $\pi_1(X,x)$, by sending $p$ to $p_*\pi_1(Y,y)$.\begin{proof}
        We have already shown this is injective, so for a subgroup $G\leq \pi_1(X,x)$ we have to construct a cover $p:(Y,y)\rightarrow (X,x)$ with $p_*\pi_1(Y,y)=G$, let $q:\tilde{X}\rightarrow X$ be the universal cover, and $H\leq \Deck(\tilde{X})$ correspond to $H$, then let $Y = \tilde{X}/H$ with basepoint $y=[\tilde{x}]$, then for all $h\in H$, we have $q(z)=q(h(z))$, the universal property gives us $p:(Y,y)\rightarrow (X,x)$ with $p\pi = q$ sending $[z]$ to $q(z)$.

        To show $p$ is a covering, choose $w\in X$ with $w\in U \subset X$ a path-connected neighbourhood with $q^{-1}(U) = \bigsqcup_{\alpha\in A} V_\alpha \subset \tilde{X}$ with each $q\vert_{V_\alpha}V_\alpha\rightarrow U$ a homeomorphism.
    \end{proof}

    TODO: finish this too
\end{theorem}

\begin{proposition}
    Any subgroup of a free group is free.
\end{proposition}

\section{Homology}

If we want an invariant that detects higher dimensional holes, the natural choice would be to consider $\pi_n(X)$ the homotopy classes of maps $S^n\rightarrow X$; unforutnately, these are very difficult to compute. By adding some combinatorial structure to our shapes, and keeping our groups abelian, we can define a functor that's much easier to compute on many of the nice spaces we care about.

\subsection{Simplicial homology}

\begin{definition}
    A \textbf{$\Delta$-complex} structure on a space $X$ is a collection of maps: \[
    \sigma_{\alpha} : \Delta^n\rightarrow X
    \]where $n$ depends on $\alpha$, with all of the $\sigma_\alpha$ satisfying the following:\begin{itemize}
        \item The restriction of each $\sigma_\alpha$ to $\mathring{\Delta}^n$ is injective;
        \item Each point $x\in X$ lies in exactly one of the $\sigma_\alpha(\mathring{\Delta}^n)$;
        \item For each face map $\delta_i$, there exists some $\sigma_\beta$ in the $\Delta$-complex with $\sigma_\alpha\circ\delta_i=\sigma_\beta$;
        \item A subset $U\subset X$ is open iff all of the $\sigma_{\alpha}^{-1}(U)$ is open in each $\Delta^n$.
    \end{itemize}
\end{definition} 

$\Delta$-complexes are a special example of CW-complexes, thus are Hausdorff. When considering these examples, it is paramount to conflate the most injective maps $\sigma_\alpha$ with their images in $X$, which will be simplices with some collapsing of their boundaries.

TODO: 

\begin{definition}
    Given a $\Delta$-complex structure on $X$, we form the \textbf{simplicial complex}: \[
    \cdots \xrightarrow{d_{n+1}^\Delta} C_n^\Delta(X) \xrightarrow{d_n^\Delta} C_{n-1}^\Delta(X) \xrightarrow{d_{n-1}^\Delta} \cdots \xrightarrow{d_3^\Delta} C_2^\Delta(X) \xrightarrow{d_2^\Delta} C_1^\Delta(X) \xrightarrow{d_1^\Delta} C_0^\Delta(X) \xrightarrow{d_0^\Delta} 0
    \] where $C_n^\Delta(X)$ is the free abelian groups on the $n$-simplices in $X$, and the differential is given by \[
    d_n^\Delta(\sigma_\alpha) = \sum_{i=0}^n(-1)^i\sigma_\alpha\circ\delta_i
    \]
\end{definition}

\begin{proposition}
    $d_\bullet^\Delta$ does in fact make $C_\bullet^\Delta$ a chain complex, i.e. $d_n^\Delta\circ d_{n+1}^\Delta=0$ for all $n$.
    \begin{proof}
        By linearity, it suffices to consider an $(n+1)$-simplex $\sigma_\alpha:\Delta^{n+1}\rightarrow X$, on which we have: \[
        d_n^\Delta d_{n+1}^\Delta(\sigma) = d_n^\Delta\pr{\sum_{i=0}^{n+1}(-1)^i\sigma\delta_i} = \sum_{i=0}^{n+1}(-1)^id_n^\Delta(\sigma\delta_i)
        = \sum_{i=0}^{n+1}(-1)^i\pr{\sum_{j=0}^{i-1}(-1)^{j}\sigma\delta_i\delta_j -
        \sum_{j=i+1}^{n+1}(-1)^{j}\sigma\delta_i\delta_j}
        \] in this sum, for each $0\leq a < b \leq n+1$, there will be two terms of $\sigma\delta_a\delta_b$, one with coefficient $(-1)^{a+b}$, and the other with $-(-1)^{a+b}$, thus the sum is zero on all $(n-1)$-simplices.
    \end{proof}
\end{proposition}



\begin{definition}
    If $X$ has a $\Delta$-complex structure, then we define the \textbf{simplicial homology} of $X$, as the homology of the chain complex $(C_\bullet^\Delta,d_\bullet^\Delta)$, we write this as $H_*^\Delta(X):= H_*(C^\Delta_\bullet(X),d)$.
\end{definition}

TODO: computations and pictures

Maps of spaces don't necessarily induce maps on chain complexes, there is no guarantee that the image of the $\Delta$-complex is a $\Delta$-complex.

TODO: give an example

\subsection{Singular homology}

We want a formulation of homology that is immediately functorial.

\begin{definition}
    To a space $X$ we can associate the \textbf{singular complex} $C_\bullet(X)$, when $C_n(X)$ is the free abelian group on maps $\sigma:\Delta^n\rightarrow X$, and the differential is as with the simplicial complex.

    From this we also define the \textbf{singular homology} of $X$, $H_*(X) := H_*(C_\bullet(X),d)$
\end{definition}

\begin{proposition}
    A continuous map $f:X\rightarrow Y$ induces homomorphisms $f_*:H_n(X)\rightarrow H_n(Y)$ for all $n$, this makes each homology group a functor $H_n(-):\textbf{Top}\rightarrow\textbf{Ab}$.
    \begin{proof}
        We will show $f$ induces a chain map $f_\sharp:C_\bullet(X)\rightarrow C_\bullet(Y)$. We define $f_\sharp$ on generators by $f_\sharp(\sigma)=f\circ\sigma$, to see this is a chain map observe: \[
            f_\sharp(d\sigma) = f_\sharp\pr{\sum_{i=0}^n(-1)^i\sigma\circ d_i} = \sum_{i=1}^n(-1)^if_\sharp(\sigma\circ d_i) = \sum_{i=1}^n(-1)^i(f\circ\sigma)\circ d_i= d(f\circ\sigma) = d(f_\sharp)
        \] All the other properties are obvious.
    \end{proof}
\end{proposition}

\begin{corollary}
    If $f:X\rightarrow Y$ is a homeomorphism, then $f_*$ is an isomorphism.
\end{corollary}

We can go further, showing homotopy equivalences induce isomorphisms on homology.

\begin{proposition}
    If $f,g:X\rightarrow Y$ are homotopic maps, then the chain maps $f_\sharp,g_\sharp$ are chain homotopic,\footnote{Two chain maps $f,g:A_\bullet\rightarrow B_\bullet$ are called \textbf{chain homotopic} if there exists some $K:A\rightarrow B[1]$ such that $f-g=dK+Kd$.} thus $f_*=g_*$.

    \begin{proof}
        Let $H:X\times I\rightarrow Y$ be the homotopy from $f$ to $g$, then $H\circ(\sigma\times\id)$ is a homotopy from $f\circ\sigma$ to $g\circ\sigma$. But the domain of this map, $\Delta^n\times I$, isn't a simplex, we can decompose it as: \[
        \Delta^n\times I = [v_0,\ldots,v_n]\times I = \bigcup_{i=0}^n[(v_0,0),\ldots,(v_i,1),(v_i,1),\ldots,(v_n,1)] =: \bigcup_{i=0}^nG_\sigma^i
        \]TODO: low dimensional picture 

        Giving us $n+1$ $n$-simplices in $Y$ associated to $\sigma$, let's defien the obvious maps $K_n:C_n(X)\rightarrow C_{n+1}(Y)$: \[
        K_n(\sigma) = \sum_{i=0}^n(-1)^iG_\sigma^i
        \] which I claim is a chain homotopy. Observe: \begin{gather*}
        dK_n(\sigma) = \sum_{i=0}^n(-1)^id(G_\sigma^i) = 
        \sum_{i=0}^n\pr{\sum_{j=0}^i (-1)^{i+j}G_\sigma^i\circ \delta_j - \sum_{j=i}^n(-1)^{i+j}G_\sigma^i\circ\delta_{j+1}}\\
        K_{n-1}(d\sigma) = \sum_{i=0}^n(-1)^iK_{n-1}(\sigma\circ\delta_i) = 
        \sum_{i=0}^n\pr{\sum_{j=0}^{i-1} (-1)^{i+j}G_\sigma^j\circ \delta_{i+1} - \sum_{j=i+1}^n(-1)^{i+j}G_\sigma^j\circ\delta_i}
        \end{gather*}
        with some careful inspection, and some picturing the $n=2$ case, we can see the only terms remaining in the sum are those where we miss a point with the index that $G$ double counts: \[
            dK_n(\sigma) + K_{n-1}d(\sigma) = \sum_{i=0}^n(G_\sigma^i\circ\delta_i - G_\sigma^i\circ\delta_{i+1}) = G_\sigma^0\circ\delta_i - G_\sigma^n\circ\delta_{n+1} = g\circ\sigma - f\circ\sigma
        \] Thus, by linearity, $K_\bullet$ is a chain homotopy.
    \end{proof}
\end{proposition}

So we can view each singular homology groups as a functor \[
H_n:\textbf{hTop}\rightarrow\textbf{Ab}
\] from the category of topological spaces and \textit{homotopy equivalence classes} of continuous maps.

\begin{corollary}
    If $f:X\rightarrow Y$ is a homotopy equivalence, then $f_*:H_n(X)\rightarrow H_n(Y)$ is an isomorphism.
\end{corollary}

\subsection{Relative homology}

\begin{lemma}
    Given a short exact sequence of chain complexes: \[
    0 \rightarrow A_\bullet \xrightarrow{i} B_\bullet \xrightarrow{j} C_\bullet \rightarrow 0
    \] there is a long exact sequence of homology groups: \[
    \cdots \rightarrow H_n(A) \xrightarrow{i_*} H_n(B) \xrightarrow{j_*} H_n(C) \xrightarrow{\partial} H_{n-1}(A) \rightarrow \cdots
    \]
    \begin{proof}
        We first define $\partial$:\[
        \begin{tikzcd}
        0 \arrow[r] & A_n \arrow[r, "i"] \arrow[d, "d"] & B_n \arrow[r, "j"] \arrow[d, "d"] & C_n \arrow[r] \arrow[d, "d"] & 0 
        & & & b \arrow[r, mapsto] \arrow[d, mapsto] & c \arrow[d, mapsto]\\
        0 \arrow[r]& A_{n-1} \arrow[r, "i"] \arrow[d, "d"] & B_{n-1} \arrow[r, "j"] \arrow[d, "d"] & C_{n-1} \arrow[r] \arrow[d, "d"] & 0 
        & & a \arrow[r,mapsto] \arrow[d, mapsto] & d(b) \arrow[r, mapsto] \arrow[d,mapsto]& 0 \\
        0 \arrow[r]& A_{n-2} \arrow[r, "i"] & B_{n-2} \arrow[r, "j"] & C_{n-2} \arrow[r]& 0 
        &  & 0 \arrow[r,mapsto] & 0\\
        \end{tikzcd}
        \] Following the above diagram: given some $[c]\in H_n(C)$, as $j$ is surjective there is some $b\in B_n$ with $j(b)=c$, by commutativity we have $j(d(b))=d(c)=0$, so $d(b)\in\ker{j}=\im{i}$, as $i$ is injective we can find a unique $a\in A_{n-1}$ with $i(a)=d(b)$, once again by commutativity $i(d(a)) = d(d(b)) = 0$, as $i$ is injective we must have $d(a) = 0$ giving a homology class $[a]\in H_n(A)$ which take as $\partial([c])$.

        We have to check our choice is well defined, as we have made a choice for $b$; if we chose a different $b'\in B_n$ with $j(b')=c$, then $j(b-b')=0$, thus $b-b'\in\im{i}$ with preimage $a'\in A_n$ so $b'=b+i(a')$, now $d(b') = d(b)+di(a') = d(b)+id(a') = i(a+d(a'))$, which will give the same homology class.

        I claim $\partial$ is now obviously a homomorphism $\ker d_{n}\rightarrow H_{n-1}(A)$.

        To show this descends to a map on homology we have to show the image of $\im d\subseteq C_n$ is $0$. If $c=d(c')$ for some $c'\in C_{n+1}$, then by surjectivity there exists some $b'\in B_{n+1}$ with $j(b')=c'$, by commutativity our $b = d(b')$ and so $i(a)=d(b)=0$, thus $a=0$.

        All of the exactness are very direct diagram chases, everything will work if you follow your nose.
    \end{proof}
\end{lemma}

\begin{corollary}
    For $A\subset X$ a subspace, write $C_\bullet(X,A)$ for $C_\bullet(X)/C_\bullet(A)$, and $H_*(X,A)$ for the homology groups of this complex, then there is a long exact sequence in \textbf{relative} homology: \[
    \cdots \rightarrow H_n(A) \xrightarrow{i_*} H_n(X) \xrightarrow{\pi_*} H_n(X,A) \xrightarrow{\partial} H_{n-1}(A) \rightarrow \cdots
    \]
\end{corollary}

\begin{corollary}
    Given a triple $A\subset B \subset X$, by the third isomorphism theorem there is a long exact sequence in relative homology: \[
    \cdots \rightarrow H_n(B,A) \xrightarrow{i_*} H_n(X,A) \xrightarrow{\pi_*} H_n(X,B) \xrightarrow{\partial} H_{n-1}(B,A) \rightarrow \cdots
    \]
\end{corollary}

\begin{lemma}[Five lemma]
    Given the following commutative diagram of abeliang groups:\[\begin{tikzcd}
        A \arrow[r, "f"] \arrow[d, "\alpha"] &
        B \arrow[r, "g"] \arrow[d, "\beta"] &
        C \arrow[r, "h"] \arrow[d, "\gamma"] &
        D \arrow[r, "k"] \arrow[d, "\delta"] &
        E \arrow[d, "\varepsilon"] \\
        A' \arrow[r, "f'", swap] &
        B' \arrow[r, "g'", swap] &
        C' \arrow[r, "h'", swap] &
        D' \arrow[r, "k'", swap] &
        E'
    \end{tikzcd}\]
    if both rows are exact and $\alpha,\beta,\delta,\varepsilon$ are all isomorphisms, then $\gamma$ is also an isomorphism.
    \begin{proof}
        We prove this with two different \textit{four lemmas}:\[
        (1)\quad
        \begin{tikzcd}
            A \arrow[r, "f"] \arrow[d, "\alpha", swap] &
            B \arrow[r, "g"] \arrow[d, "\beta", swap] &
            C \arrow[r, "h"] \arrow[d, "\gamma"] &
            D \arrow[d, "\delta"] \\
            A' \arrow[r, "f'", swap] &
            B' \arrow[r, "g'", swap] &
            C' \arrow[r, "h'", swap] &
            D'
        \end{tikzcd}
        \qquad \qquad
        (2)\quad
        \begin{tikzcd}
            B \arrow[r, "g"] \arrow[d, "\beta", swap] &
            C \arrow[r, "h"] \arrow[d, "\gamma", swap] &
            D \arrow[r, "k"] \arrow[d, "\delta"] &
            E \arrow[d, "\varepsilon"] \\
            B' \arrow[r, "g'", swap] &
            C' \arrow[r, "h'", swap] &
            D' \arrow[r, "k'", swap] &
            E'
        \end{tikzcd}
        \]
    For $(1)$ we assume both rows are exact, $\beta,\delta$ are injective, and $\alpha$ is surjectve; and show $\gamma$ is injective:

        Have $c\in C$ such that $\gamma(c)=0\in C'$ so $\delta h(c) = h'\gamma(c) = 0$ and as $\delta$ is injective, $h(c)=0$;

        so $c\in \ker h = \im g$, there exists $b\in B$ with $c=g(b)$ and $g'(\beta(b)) = \gamma(g(b)) = \gamma(c) = 0$;

        so $\beta(b)\in\ker g'=\im f'$, there exists $a'\in A'$ with $f'(a') = \beta(b)$;

        $\alpha$ surjective $\implies$ there exists $a\in A$ with $a' = \alpha(a)$ and $\beta(b) = f'(\alpha(a)) = \beta(f(a))$;

        $\beta$ injective $\implies b=f(a)$ so $c=g(b) = f(g(a)) = 0$.

    \noindent For $(2)$ we assume both rows are exact, $\beta,\delta$ are surjective, and $\varepsilon$ is injective; and show $\gamma$ is surjective:
        
        Have $c'\in C'$; as $\delta$ is surjective there is a $d\in D$ with $\delta(d)=h'(c')$;
        
        $\varepsilon (k(d)) = k'\delta(d) = k'h'(c) = 0$; as $\varepsilon$ is injective, $k(d) = 0$;
        
        so $d\in\ker k = \im h$, there exists $c\in C$ with  $h(c) = d$; so $h'(\gamma(c))\delta(h(c)) = \delta(b) = c'$;
        
        so $c'-\gamma(c)\in\ker h' = \im g'$, there is $b'\in B'$ with $g'(b')=c'-\gamma(c)$;
        
        as $\beta$ is surjective, there is $b\in B$ with $\beta(b) = b'$; so $\gamma(g(b)) = g'(\beta(b)) = c'-\gamma(c)$;
        
        so $c'=\gamma(g(b))+\gamma(c) = \gamma(g(b) + c)$.

    \noindent Therefore $\gamma$ is both injective and surjective, thus an isomorphism.
    \end{proof}
\end{lemma}

All of the previously established properties of homology also hold for relative homology as a functor from the category of pairs of spaces and continuous functions preservings subspaces to \textbf{Ab}.

\begin{corollary}
    Let $f:(X,A)\rightarrow (Y,B)$ be a map of pairs. If any tow of the follows maps are isomorphisms, the third is as well: \[
    (f\vert_A)_* : H_*(A)\rightarrow H_*(B) \qquad f_*:H_*(X)\rightarrow H_*(Y) \qquad f_*: H_*(X,A) \rightarrow H_*(X,B)
    \]\begin{proof}
        Apply the five lemma to the induced maps between the long exact sequences of pairs.
    \end{proof}
\end{corollary}

\subsection{Barycentric subdivision}

\begin{definition}
    Let $\cU=\{U_i\}$ be a collection of subsets of $X$ whose interiors form an open cover of $X$. For each $n$ consider \[
    C^\cU_n(X) = \sum_iC_n(U_i) \subseteq C_n(X)
    \] equipping this with the usual differential gives us the chain complex $C_\bullet^\cU(X)$ with the canonical natural map $i:C^\cU_\bullet(X)\rightarrow C_\bullet(X)$ which induces $i_*:H_n^\cU(X)\rightarrow H_n(X)$.
\end{definition}

With this construction, we want to approximate a general homology class as a sum of small simplices, each of which lie entirely within one component of an open cover. This will be reified with an isomorphism between the $H_*(X)$ and $H_*^\cU(X)$ above, and our exact method for replacing our simplices with smaller ones will be \textit{barycentric subdivision}.

We define the barycentric subdivision of the $n$-simplex $\Delta^n=[v_0,\ldots,v_n]\subset \bR^{n+1}$ inductively. Let $b=\frac{1}{n+1}(v_0+\cdots+v_n)$, then for each facet $[w_1,\ldots,w_n]$ of $\Delta^n$, consider the $n$-simplex $[b,w_1,\ldots,w_n]$. This collection of $(n+1)!$ sipmlices are the barycentric subdivision of $\Delta^n$.

TODO: some pictures in multiple dimensions of multiple iterations of barycentric subdivision.

To show these simplices get appropriately small, recall this lemma from analysis.

\begin{proposition}[Lebesgue covering lemma]
    Let $\{U_i\}$ be an open cover of a compact metric space $X$, there exists some $\delta>0$ such that for all $x\in X$, the ball $B_\delta(x)$ is wholy contained in one of the $U_i$.
    \begin{proof}
        As $X$ is compact, wlog we may take $\{U_i\}$ to be finite. So we can consider the function $f:X\rightarrow \bR$ given by \[
        f(x) := \frac{1}{n}\sum_{i=0}^nd(x,U_i^c)
        \] as $f$ is continuous and defined on a compact set, it must achieve a minimum $\delta$. And, as every $x$ is contained in some $U_i$, $f$ is never $0$, thus $\delta>0$. Now, for all $x\in X$, there exists some $U_i$ such that $d(x,U_i^c)\geq\delta$, thus $B_\delta(x)\subseteq U_i$.
    \end{proof}
\end{proposition}

\begin{lemma}
    Let $\{U_i\}$ be an open cover of the $n$-simpliex $\Delta^n$, then we can barycentrically subdivide $\Delta$ finitely many times so that every simpliex is wholy contained in one of the $U_i$.
    \begin{proof}
        Let $\delta>0$ be the Lebesgue number for this open cover, then every simplex of diameter $<\delta$ will lie whole in one $U_i$. So it suffices to show the diameter of the subdivided simplices can be made arbitrarily small with a finite number of subdivisions.

        Suppose $\Delta'$ is a simplex in the barucentric subdivision of $\Delta$, then I claim: \[
        \diam(\Delta')\leq \frac{n}{n+1}\diam(\Delta)
        \] which is clearly sufficient. The diameter of a linear simplex in $\bR^n$ is in fact equal to the distance between a pair of vertices, so for all pairs of vertices $x,y\in\Delta'$ we just have to consider two cases:\begin{itemize}
            \item If neither $x$ nor $y$ is the barycenter $b$, then they belong to a simplex in the barycentric subdivision of a facet of $\tilde{\Delta}\subset \Delta$, so by induction \[
            \abs{x-y} \leq \frac{n-1}{n}\diam(\tilde{\Delta}) \leq \frac{n}{n+1}\diam(\tilde{\Delta})\leq \frac{n}{n+1}\diam(\Delta)
            \]
            \item If instead $y$ is the barycenter $b$, then if we let $b_0$ be the barycenter of the facet of $\Delta$ opposite $x=v_0$, we have \vspace{-12pt}\begin{gather*}
                \abs{x-y} = \abs{v_0 - \frac{v_0+\cdots+v_n}{n+1}}
                = \abs{\frac{n}{n+1}v_0 - \frac{n}{n+1}\pr{\frac{v_1+\cdots v_n}{n}}} = \\\frac{n}{n+1}\abs{v_0-b_0}\leq \frac{n}{n+1}\diam(\Delta)\qedhere
            \end{gather*}
        \end{itemize}
    \end{proof}
\end{lemma}
We now just have to put barycentric subdivision into the form of a chain map. We will do this inductively: for a singular $n$-simplex $\sigma:\Delta^n\rightarrow X$ consider each of its facets $\sigma\circ\delta_i$, and consider its barycentric subdivision consisting of $(n-1)$-simplices. Consider the cone of each of these simplices, and send $\sigma$ to the alternating sums of the subdivision of facets.\[
S(\sigma) = c_b(S(d \sigma))
\]

\end{document}